\section{Profiles with large mean gradient}
\label{sec:profiles}

A smooth real function $\phi$ will determine the amplitude envelope
$e^{n\phi}$ at frequency $n$ for the waves in Section~\ref{sec:waves}.
Their construction
requires a strict Hessian inequality, whereas their eventual nodal-area
lower bound grows with the integrated gradient of $\phi$.  We show that
these two requirements are compatible: a local perturbation can increase
the gradient integral arbitrarily while preserving the Hessian condition
and changing $\phi$ arbitrarily little in $C^0$.  We then apply this
construction on $\Sph^3$, preserving a prescribed radial profile
outside one coordinate cube.

Throughout this section, gradients, Hessians, and norms are computed using a
fixed smooth Riemannian metric $g$ on a three-dimensional manifold.  Write
\[
 a_\phi=\grad_g\phi,\qquad
 L_\phi=(1+|a_\phi|^2)^{1/2},\qquad H_\phi=\Hess_g\phi.
\]

\begin{definition}[Strict and full profiles]
\label{def:strict-profile}
A smooth function $\phi$ is \emph{strict} at a point if there exists a real
vector $b$ at that point such that
\begin{equation}
 b\perp a_\phi,\qquad |b|=L_\phi,\qquad
 H_\phi(a_\phi,a_\phi)+H_\phi(b,b)>0.
 \label{eq:profile-strict}
\end{equation}
It is \emph{full} at that point if the last inequality holds for every $b$
satisfying the first two conditions.
\end{definition}

For a symmetric form $H$ and a unit vector $v\perp a$, it is convenient to set
\[
 \mathcal T(H,a;v)=H(a,a)+(1+|a|^2)H(v,v).
\]
Fullness on a compact set where $a_\phi$ does not vanish gives a positive
minimum of $\mathcal T(H_\phi,a_\phi;v)$ over the transverse unit-circle
bundle.  That minimum is allowed to depend on the profile.

The main construction in this section is local.

\begin{proposition}[Amplification of the mean gradient]
\label{prop:profile-amplification}
Let $\Omega$ be a coordinate cube with compact closure in a smooth
three-dimensional Riemannian manifold.  Suppose that $\phi$ is smooth on a
neighborhood of $\overline\Omega$, is strict there, and satisfies
$a_\phi\ne0$ there.  Given $T>0$ and $\varepsilon>0$, there is a smooth
function $\widehat\phi$ on the same neighborhood such that
\[
 \supp(\widehat\phi-\phi)\Subset\Omega,\qquad
 \|\widehat\phi-\phi\|_{C^0(\Omega)}<\varepsilon,
 \qquad
 \int_\Omega |a_{\widehat\phi}|\,\dd V_g>T.
\]
The function $\widehat\phi$ is strict and has nonvanishing gradient on
$\overline\Omega$.

More precisely, there is a constant $\eta>0$, determined only by the fixed
periodic functions chosen below, such that one finite cycle of the
construction multiplies
$\int_\Omega|a_\phi|\,\dd V_g$ by at least $1+\eta$.  The cycle can have
arbitrarily small $C^0$ size.  Its patch sizes and frequency thresholds may
depend on the entire incoming profile and metric.
\end{proposition}

The cycle has two parts.  A first perturbation makes the profile full on a
fixed positive fraction of its gradient mass while changing the gradient
arbitrarily little.  A radial perturbation in the two transverse variables
then increases the gradient on a fixed fraction of that region.  The radial
perturbation preserves strictness, allowing the cycle to be repeated.

\subsection{Localization and periodic averaging}

We use two elementary facts in both parts of the construction.

\begin{lemma}[Disjoint patches and averaging]
\label{lem:profile-localization}
Let $U$ be an open subset of a relatively compact coordinate cube, and let
$\mu$ be a finite measure given by a positive smooth density.
\begin{enumerate}
 \item Given an open neighborhood of each point of $U$, one can choose
 finitely many coordinate cubes $P_j$, with disjoint interiors and closures
 contained in those neighborhoods and in $U$, and functions
 $0\le\chi_j\le1$ in $C_c^\infty(P_j)$, such that $\chi_j=1$ on open
 cores $C_j\Subset P_j$ and
 \[
   \mu\Bigl(\bigcup_j C_j\Bigr)\ge\frac12\mu(U).
 \]
 In particular, the supports of the $\chi_j$ are pairwise disjoint.
 \item Suppose that $y$ is a smooth submersion to $\R^q$, with $q=1$ or
 $2$, on a neighborhood of $\overline P$, and that $C\Subset P$ is
 measurable.  Let $E\subset\R^q$ be periodic with respect to a full
 lattice, have boundary of measure zero in a period cell, and occupy a
 fraction $p$ of that cell.  Then
 \[
  \lim_{N\to\infty}
  \mu\{x\in C:Ny(x)\in E\}=p\,\mu(C).
 \]
\end{enumerate}
\end{lemma}

\begin{proof}
For the first assertion, choose a compact subset of $U$ carrying more than
three quarters of its mass.  A finite subcover of the prescribed
neighborhoods has a Lebesgue number on this compact set.  A sufficiently fine
coordinate grid therefore gives finitely many cubes, each with closure in
one of the prescribed neighborhoods and in $U$, whose union contains that
compact set except for grid faces.  The faces have zero mass.  Shrinking the
cores inside these cubes loses arbitrarily little mass, and smooth cutoffs
supported in the cubes can be chosen to equal one on the cores.

For the second assertion, submersion coordinates and a finite partition of
unity express the pushforward of $\mu|_C$ under $y$ as an integrable density
on $\R^q$.  Thus it suffices to average the bounded periodic function
$\mathbf 1_E(N\cdot)$ against an $L^1$ function.  For the indicator of a
rectangle, the assertion follows by tiling with period cells and estimating
the cells meeting the rectangle's boundary.  Approximation of an $L^1$
function by finite linear combinations of rectangle indicators proves the
claim.  The boundary assumption on $E$ ensures that its cell average is
well defined and unchanged by choices on its boundary.
\end{proof}

We fix a smooth function $h$ of period one such that
\[
 h''(t)=2\cos(2\pi t).
\]
For example, $h(t)=-(2\pi^2)^{-1}\cos(2\pi t)$.  The periodic set
\[
 E_h=\{t:h''(t)>1\}
\]
has fraction $p_h=1/3$ in each period.

\subsection{Preparing a full region}

\begin{lemma}[Preparation]
\label{lem:profile-preparation}
Under the hypotheses of Proposition~\ref{prop:profile-amplification}, let
$0<\epsilon_p\le1/3$ and $\varepsilon>0$.  There is a smooth strict
function $\psi$, with nonvanishing gradient on $\overline\Omega$, such that
\[
 \supp(\psi-\phi)\Subset\Omega,\qquad
 \|\psi-\phi\|_{C^0(\Omega)}<\varepsilon,
 \qquad |a_\psi-a_\phi|\le\epsilon_p|a_\phi|
 \quad\hbox{on }\overline\Omega.
\]
Its full region $U\subset\Omega$ satisfies
\begin{equation}
 \int_U|a_\psi|\,\dd V_g
 \ge\frac{p_h}{8}\int_\Omega|a_\psi|\,\dd V_g.
 \label{eq:prepared-mass}
\end{equation}
The change $\psi-\phi$ can, in addition, be made arbitrarily small in
$C^1$.
\end{lemma}

\begin{proof}
Put $a=a_\phi$.  At each point, choose a strict unit direction $v\perp a$.
Because $a$ does not vanish, the flow-box construction gives a local first
integral $s$, with $\dd s(a)=0$, whose differential at the point is any
prescribed nonzero covector annihilating $a$.  Prescribe it also to
annihilate $v$.  On a sufficiently small neighborhood, the line
\[
 a^\perp\cap\ker\dd s
\]
continues to be strict, with a positive margin on a smaller compact
neighborhood.  The differential $\dd s$ remains nonzero there.

Apply Lemma~\ref{lem:profile-localization} to these neighborhoods with
$\dd\mu=|a|\,\dd V_g$.  It gives finitely many cubes and cutoff-one cores
capturing at least half the mass.  We first describe the perturbation on
one cube, suppressing its index.  Let $\chi$ be its cutoff.  For a constant
$C>0$ to be fixed before $N$, set
\[
 \psi_N=\phi+C N^{-2}\chi h(Ns).
\]
Direct differentiation gives, uniformly on the patch,
\begin{align*}
 a_{\psi_N}
 &=a+\frac C N\chi h'(Ns)\grad s
       +\frac C{N^2}h(Ns)\grad\chi,\\
 H_{\psi_N}
 &=H_\phi+C\chi h''(Ns)\,\dd s\otimes\dd s+O(N^{-1}).
\end{align*}
All constants in this proof may depend on the fixed patch data, including
$C$, but not on $N$.

For large $N$, choose a unit vector $v_N$ on the line
\[
 a_{\psi_N}^{\perp}\cap\ker\dd s.
\]
This line tends uniformly to the old strict line.  The added leading
Hessian has zero value on $v_N$, while
$\dd s(a_{\psi_N})=O(N^{-1})$.  Its value on the new gradient is
therefore $O(N^{-2})$, even on the portions where $h''$ is negative.  The
remaining Hessian error is $O(N^{-1})$.  The old strict margin consequently
preserves strictness for all sufficiently large $N$.

We choose $C$ to obtain fullness on the cutoff-one core where $Ns\in E_h$.
For the old gradient, consider the symmetric form on $a^\perp$
\[
 B=(1+|a|^2)H_\phi|_{a^\perp}+H_\phi(a,a)\Id.
\]
In an orthonormal basis $(v,w)$ with $\dd s(v)=0$ and
$w=\grad s/|\dd s|$, its $vv$ entry is bounded below by some $m>0$.
All its entries are bounded on the patch.  Adding
\[
 C(1+|a|^2)|\dd s|^2\,w\otimes w
\]
makes $B$ uniformly positive definite when $C$ is sufficiently large.
Indeed, if $|B_{vw}|\le M$ and $B_{ww}\ge-M$, it is enough to make the
new $ww$ entry exceed $M^2/m+m$; the Schur complement is then positive.
The factor $(1+|a|^2)|\dd s|^2$ has a positive patch minimum, so a finite
$C$ suffices.  On the specified core $\chi=1$ and $h''>1$, so the actual
leading addition is at least this large.  Changing the gradient by
$O(N^{-1})$ and adding the Hessian remainder preserves its positive
minimum for large $N$.  Thus $\psi_N$ is full there.

Perform these perturbations on the finite disjoint patches, with their
chosen constants $C$, using a common sufficiently large $N$.  The changes
are arbitrarily small in $C^1$, preserve nonvanishing, and have compact
support in $\Omega$.  Lemma~\ref{lem:profile-localization} shows that the
union $G_N$ of the cutoff-one regions with $Ns\in E_h$ satisfies
\[
 \int_{G_N}|a|\,\dd V_g
 \ge\frac{p_h}{4}\int_\Omega|a|\,\dd V_g
\]
for sufficiently large $N$.  Choose $N$ still larger so that the stated
pointwise relative gradient bound holds.  With
$\mathcal M(\zeta)=\int_\Omega|a_\zeta|\,\dd V_g$, the full region
therefore has mass at least
\[
 (1-\epsilon_p)\frac{p_h}{4}\mathcal M(\phi)
 \ge\frac{p_h(1-\epsilon_p)}{4(1+\epsilon_p)}\mathcal M(\psi_N)
 \ge\frac{p_h}{8}\mathcal M(\psi_N).
\]
Fullness is open where the gradient is nonzero, so this gives the asserted
open full region.  The $C^0$ bound and any additional prescribed $C^1$
smallness follow by increasing $N$ once more.
\end{proof}

\subsection{Radial corrugation in two transverse variables}

Fix numbers
\[
 0<r_0<r_1<r_2<R<\frac14
\]
and a smooth nonnegative function $f'$ supported in $(r_0,R)$, with
$f'\ge c_1>0$ on $(r_1,r_2)$.  Define
$f(s)=\int_0^s f'(t)\,\dd t$.  On every disk of radius $R$ centered at
a point of $\mathbb Z^2$, define $F$ by the radial function $f$; extend it
by the constant $f(R)$ off these disks.  The result is a smooth
$\mathbb Z^2$-periodic function, constant near the disk centers and
boundaries.  Its annular set
\[
 E_F=\bigcup_{k\in\mathbb Z^2}
       \{y:r_1<|y-k|<r_2\}
\]
has a fixed positive fraction $p_F=\pi(r_2^2-r_1^2)$ in a period cell.

\begin{lemma}[Transverse radial corrugation]
\label{lem:radial-corrugation}
There are constants $\delta>0$ and $g_0>0$, depending only on the fixed
function $F$, with the following property.  Let $\psi$ be full with
nonvanishing gradient on a neighborhood of the closure of a relatively
compact coordinate patch $P$.  Suppose that
$y=(y_1,y_2)$ is a smooth submersion there, with
\[
 \dd y(a_\psi)=0,
 \qquad\text{both nonzero singular values of }\dd y
       \text{ in }[1/2,2].
\]
Let $A>0$ be constant and satisfy $1/2\le A/|a_\psi|\le2$, and let
$0\le\chi\le1$ belong to $C_c^\infty(P)$.  Set
\[
 \Phi_N=\psi+\frac{\delta A}{N}\chi F(Ny).
\]
For every sufficiently large integer $N$, the function $\Phi_N$ is strict
and has nonvanishing gradient on $\overline P$.  Moreover,
$\|\Phi_N-\psi\|_{C^0(P)}\to0$.  Given any $\epsilon_r>0$, all
sufficiently large $N$ satisfy
\begin{align}
 |a_{\Phi_N}|&\ge(1-\epsilon_r)|a_\psi|
                      &&\text{on }\overline P,\label{eq:radial-no-loss}\\
 |a_{\Phi_N}|&\ge(1+g_0-\epsilon_r)|a_\psi|
                      &&\text{where }\chi=1\text{ and }Ny\in E_F.
                      \label{eq:radial-gain}
\end{align}
The frequency thresholds may depend on all the fixed patch data and on
$\epsilon_r$; the constants $\delta,g_0$ do not.
\end{lemma}

\begin{proof}
Put $a=a_\psi$ and $a_N=a_{\Phi_N}$.  When $Ny$ lies in a periodic disk,
let $d_R,d_T$ be the radial and angular unit-coefficient combinations of
$\dd y_1,\dd y_2$.  The singular-value hypothesis gives
\[
 \frac12\le |d_R|,|d_T|\le2.
\]
These covectors need not be orthogonal in the metric.  Outside the disks,
set the leading derivatives $f',f'',f'/s$ to zero.  There, or in the
constant central portions of the disks, choose any orthonormal coefficient
pair when such notation is needed.  Set $f'/s=0$ at the center.  Since $f$
is constant near the center, this agrees with its smooth continuation.

The exact gradient formula and the Hessian expansion are
\begin{equation}
 \begin{aligned}
 a_N&=a+q\,d_R^\sharp+e_N,\\
 q&=\delta A\chi f',\qquad
 e_N=\frac{\delta A}{N}F(Ny)\grad\chi,
 \end{aligned}
 \label{eq:radial-gradient}
\end{equation}
\begin{equation}
 H_{\Phi_N}
 =H_\psi+N\delta A\chi
       \left[f''d_R^2+\frac{f'}s d_T^2\right]+E_N,
 \label{eq:radial-hessian}
\end{equation}
with
\begin{equation}
 |E_N|\le C_P\left((\chi+|\dd\chi|)f'+N^{-1}\right).
 \label{eq:radial-remainder}
\end{equation}
Here and below $C_P$ denotes a finite constant depending on the fixed
patch data, independent of $N$.  To check the remainder, its terms are
\begin{align*}
 &\delta A\chi\sum_i F_i(Ny)\Hess y_i,\\
 &\delta A\left(\dd\chi\otimes\sum_i F_i(Ny)\dd y_i
       +\sum_i F_i(Ny)\dd y_i\otimes\dd\chi\right),\\
 &\frac{\delta A}{N}F(Ny)\Hess\chi.
\end{align*}
Since $|(F_1,F_2)|=f'$, these give
Equation~\eqref{eq:radial-remainder}.  In the constant portions of $F$,
its value contributes only to $e_N$ and to the last displayed term.

We first control the leading Hessian trace in
\eqref{eq:radial-hessian}.  The leading increment in
\eqref{eq:radial-gradient} is perpendicular to $a$, so
\[
 |a_N|\ge |a|-|e_N|.
\]
Thus $a_N$ is nonzero for large $N$.  Choose a vector $b_N$ satisfying
\[
 b_N\perp a_N,\qquad d_R(b_N)=0,\qquad
 |b_N|=L_N:=(1+|a_N|^2)^{1/2}.
\]
The two linear constraints have a nonzero common solution in dimension
three.  They imply
\[
 \langle a,b_N\rangle=-\langle e_N,b_N\rangle.
\]
Because $|a|$ has a positive patch minimum, the component of $b_N/L_N$
along $a$ tends uniformly to zero.  The kernel of $\dd y$ is precisely
the line spanned by $a$.  Its singular-value lower bound and the identity
$d_R(b_N)=0$ consequently give, for all sufficiently large $N$,
\[
 |d_T(b_N)|=|\dd y(b_N)|\ge\frac14 L_N.
\]
Also $A\le3L_N$ for large $N$, and
\[
 |d_R(a_N)|\le4\delta A\chi f'+2|e_N|.
\]
With $K=\|f''\|_\infty$ and
$B_N=f''d_R^2+(f'/s)d_T^2$, it follows that
\begin{align*}
 B_N(a_N,a_N)+B_N(b_N,b_N)
 &\ge\frac1{16}L_N^2\frac{f'}s-K|d_R(a_N)|^2\\
 &\ge\frac1{16}L_N^2\frac{f'}s
       -288K\delta^2L_N^2(f')^2-8K|e_N|^2.
\end{align*}
The estimate uses $(x+y)^2\le2x^2+2y^2$; in particular, its final error is
$O_P(N^{-2})$.  Write $S_0=\sup_s sf'(s)>0$.  Choose once and for all
\[
 0<\delta^2\le\frac1{9216KS_0}.
\]
Then $(f')^2\le S_0 f'/s$ gives
\begin{equation}
 \begin{split}
 N\delta A\chi
   \bigl[B_N(a_N,a_N)+B_N(b_N,b_N)\bigr]
 \ge N\delta A\chi\frac{L_N^2}{32}\frac{f'}s
       -C_P N^{-1}.
 \end{split}
 \label{eq:radial-leading-trace}
\end{equation}
All constants used to choose $\delta$ depend only on $F$ and the
numerical singular-value bounds.

We next control the transition where the favorable leading term can be
small.  A nonnegative smooth cutoff satisfies
\[
 |\dd\chi|\le C_P\sqrt\chi.
\]
Indeed, extend it by zero in the coordinate chart.  If its Euclidean
Hessian is bounded by $M>0$, the Taylor inequality at
$x-M^{-1}D\chi(x)$ gives
$|D\chi(x)|^2\le2M\chi(x)$; metric equivalence gives the stated bound.
Consequently, with $t=\chi f'$, Equations~\eqref{eq:radial-gradient}
and~\eqref{eq:radial-remainder} imply
\begin{equation}
 |a_N-a|\le C_P(t+N^{-1}),\qquad
 |E_N|\le C_P(t+\sqrt t+N^{-1}).
 \label{eq:radial-small-product}
\end{equation}

Let $m>0$ be the minimum of the full trace
$\mathcal T(H_\psi,a;v)$ over the transverse unit circles on
$\overline P$.  Project $b_N/L_N$ onto $a^\perp$ and normalize; for
large $N$ this gives a unit vector $v_N\perp a$ differing from
$b_N/L_N$ by $O_P(N^{-1})$.  Using also
$|L_N-L_\psi|\le|a_N-a|$, the base Hessian trace is at least
\[
 H_\psi(a_N,a_N)+H_\psi(b_N,b_N)
 \ge m-C_P(t+N^{-1}).
\]
The vectors $a_N,b_N$ are uniformly bounded for this fixed patch, so the
trace of the remainder in \eqref{eq:radial-small-product} is
bounded in absolute value by $C_P(t+\sqrt t+N^{-1})$.
Choose $t_0>0$ sufficiently small that these terms involving $t$ and
$\sqrt t$ total less than $m/4$ when $t\le t_0$.  For sufficiently large
$N$, Equation~\eqref{eq:radial-leading-trace} now gives a strictly
positive total trace throughout $t\le t_0$.

On the complementary region $t\ge t_0$, one has $s\le R$, and the
positive term in \eqref{eq:radial-leading-trace} is at least
\[
 \frac{N\delta A t_0}{32R}.
\]
It tends to infinity and dominates the bounded base and remainder traces.
This proves strictness for all sufficiently large $N$.  In this argument
$t_0$ and the frequency threshold can depend on the full margin $m$ and
all patch derivatives.  The choice of $\delta$ is already fixed.  In
particular, the region $t\to0$ requires no lower bound for the product
$Nt$.

Finally, exact orthogonality in the leading gradient gives
\[
 |a+q d_R^\sharp|^2=|a|^2+q^2|d_R|^2.
\]
The uniform error $e_N=O_P(N^{-1})$ proves
\eqref{eq:radial-no-loss}.  Where $\chi=1$ and $Ny\in E_F$,
the bounds $A\ge|a|/2$ and $|d_R|\ge1/2$ give
\[
 |a+q d_R^\sharp|
 \ge |a|\sqrt{1+\frac{\delta^2c_1^2}{16}}.
\]
Set
\[
 g_0=\sqrt{1+\frac{\delta^2c_1^2}{16}}-1>0.
\]
Decreasing the fixed $\delta$ if necessary, assume $g_0\le1$.  The
positive minimum of $|a|$ allows $|e_N|\le\epsilon_r|a|$ for large
$N$, proving \eqref{eq:radial-gain}.  The $C^0$ convergence
follows directly from the factor $N^{-1}$ in the definition of $\Phi_N$.
\end{proof}

\subsection{A fixed gain and finite iteration}

\begin{proof}[Proof of Proposition~\ref{prop:profile-amplification}]
Keep the periodic functions and the constants from
Lemma~\ref{lem:radial-corrugation} fixed.  Define
\[
 \theta=\frac{p_h p_F}{64},\qquad
 t_* =\theta g_0,\qquad \eta=\frac{t_*}{4}>0.
\]
We prove the cycle statement, allowing any prescribed positive $C^0$
budget for that cycle.

Apply Lemma~\ref{lem:profile-preparation} with
$\epsilon_p\le\min(1/3,t_*/16)$ and half the $C^0$ budget.  Write
$\psi$ for the resulting profile and $U$ for its full region.  The
preparation, including its frequency, is now held fixed.  On a sufficiently
small neighborhood of each point of $U$, the flow-box construction gives
two first integrals $y_1,y_2$ of $a_\psi$.  Their differentials can be
prescribed to form an orthonormal basis of $a_\psi^\perp$ at the center;
after shrinking, the nonzero singular values of $\dd y$ lie in $[1/2,2]$.
One can also choose a constant $A$ on that patch with
$1/2\le A/|a_\psi|\le2$.

Apply Lemma~\ref{lem:profile-localization} in $U$ with density
$|a_\psi|\,\dd V_g$.  Its radial patch cores capture at least half the
full-region mass.  In view of Equation~\eqref{eq:prepared-mass}, their
total mass is at least $(p_h/16)\mathcal M(\psi)$.  For all sufficiently
large common radial frequencies $N$, periodic averaging gives a set of
cutoff-one annular points of mass at least
\[
 \frac{p_h p_F}{32}\mathcal M(\psi)
 \ge\theta\mathcal M(\psi).
\]
Use Lemma~\ref{lem:radial-corrugation} on these disjoint patches.  Increase
$N$ to obtain strictness, nonvanishing, the remaining half of the $C^0$
budget, and relative errors
\[
 \epsilon_r\le\frac{t_*}{8},\qquad \epsilon_r\le\frac{g_0}{2}.
\]
Call the resulting profile $\widehat\psi$.  On the selected annular set
the gradient multiplier is at least $1+g_0/2$; everywhere else it is at
least $1-\epsilon_r$.  Therefore
\begin{align*}
 \mathcal M(\widehat\psi)
 &\ge\left(1-\epsilon_r+\frac{\theta g_0}{2}\right)
          \mathcal M(\psi)\\
 &\ge\left(1+\frac{3t_*}{8}\right)\mathcal M(\psi)\\
 &\ge\left(1-\frac{t_*}{16}\right)
          \left(1+\frac{3t_*}{8}\right)\mathcal M(\phi)\\
 &\ge(1+\eta)\mathcal M(\phi).
\end{align*}
The last inequality uses $0<t_*\le1$.  This establishes a gain independent
of the incoming profile and of the full-region patch sizes.  The latter
are fixed only after preparation; their finite bounds determine how large
the radial frequency must be.

Since $\mathcal M(\phi)>0$, finitely many cycles raise it above $T$.
Divide $\varepsilon$ among those cycles.  Each cycle has compact support
in $\Omega$ and preserves strictness and nonvanishing on
$\overline\Omega$.  Their finite union of supports is compactly contained
in $\Omega$, and their total $C^0$ change is less than $\varepsilon$.
All functions produced are smooth.  This completes the proof of
Proposition~\ref{prop:profile-amplification}.
\end{proof}

\subsection{Profiles on the round sphere}

We record the global profiles used below, including the behavior at their
critical points.  On the unit sphere $\Sph^3\subset\R^4$, let $g_*$ be
the round metric and set
\[
 r=(x_1^2+x_2^2)^{1/2}.
\]
Fix a coordinate cube $\Omega$ whose closure lies in
$1/20<r<1/8$, and fix its coordinate Lebesgue measure $\dd x$.

\begin{proposition}[Spherical profiles]
\label{prop:spherical-profiles}
There is a $C^2$ neighborhood $\mathcal U$ of $g_*$ such that the
following holds.  For every smooth $g\in\mathcal U$, every
$1/5\le t\le3/10$, and every $T,\varepsilon>0$, there is a smooth real
function $\phi$ on $\Sph^3$ with
\[
 \phi_0=\log t+r^2-t^2,\qquad
 \supp(\phi-\phi_0)\Subset\Omega,
 \qquad \|\phi-\phi_0\|_{C^0}<\varepsilon,
\]
such that
\begin{enumerate}
 \item $\phi$ is strict on $r\le3/10$;
 \item its critical points in this region are exactly $r=0$, and at
 each such point $H_\phi$ is positive on a two-plane;
 \item $\phi>\log r$ for $0<r<t$, while $\phi<\log r$ for
 $t<r\le1$;
 \item $\displaystyle\int_\Omega |a_\phi|\,\dd V_g>T$.
\end{enumerate}
The metric-coordinate comparison constants on $\Omega$, including the
Lipschitz constants of its two-coordinate projections for the ambient
Riemannian distance, can be chosen uniformly for $g\in\mathcal U$.
In particular,
\begin{equation}
 \int_\Omega L_\phi\,\dd x
 \ge c\int_\Omega |a_\phi|\,\dd V_g
 \label{eq:profile-coordinate-mass}
\end{equation}
with one constant $c>0$ for this neighborhood.  Each constructed profile
is fixed before a wave frequency is chosen.
\end{proposition}

\begin{proof}
The function $\phi_0$ is smooth on the sphere.  Its round Hessian is
\begin{equation}
 \Hess_{g_*}\phi_0
 =2(\dd x_1^2+\dd x_2^2)-2r^2g_*.
 \label{eq:round-profile-hessian}
\end{equation}
The eigenvalues of $\dd x_1^2+\dd x_2^2$ on the tangent three-plane are
$1$, $1-r^2$, and $0$: its two nonzero eigenvalues are those of the Gram
matrix $\Id_2-(x_1,x_2)^T(x_1,x_2)$.  The Hessian eigenvalues in
\eqref{eq:round-profile-hessian} are consequently
\[
 2-2r^2,\qquad 2-4r^2,\qquad -2r^2.
\]
On $r\le3/10$, the first two exceed $1$, and the third exceeds $-1/2$.
These inequalities persist uniformly in a sufficiently small $C^2$
neighborhood $\mathcal U$ of the round metric.  The derivatives of
$\phi_0$ do not depend on $t$, so the same neighborhood works for all
the stated $t$.

For such a metric, intersect a positive two-plane of $H_{\phi_0}$ with
$a_{\phi_0}^\perp$ and choose a unit vector $v$ in the intersection.
Then
\[
 \mathcal T(H_{\phi_0},a_{\phi_0};v)
 \ge-\frac12|a_{\phi_0}|^2+1+|a_{\phi_0}|^2>0.
\]
Thus $\phi_0$ is strict throughout $r\le3/10$.  Also
\[
 |\grad_{g_*}(r^2)|^2=4r^2(1-r^2).
\]
Vanishing of a differential is independent of the metric, so $\phi_0$
has nonvanishing gradient for $0<r<1$, in particular on a neighborhood of
$\overline\Omega$.

For the envelope comparison, put
\[
 j_t(r)=\log t+r^2-t^2-\log r.
\]
One has $j_t(t)=0$ and $j_t'(r)=2r-r^{-1}$.  Hence $j_t(r)>0$ for
$0<r<t$.  On $t<r\le1$, the only possible interior critical point is a
minimum, while
$j_t(1)=\log t+1-t^2<0$ for the stated range of $t$; thus $j_t<0$ there.
Compactness gives a positive lower bound for $j_t$ on
$\overline\Omega$, uniform over $t\in[1/5,3/10]$.

Apply Proposition~\ref{prop:profile-amplification} to $\phi_0$ on
$\Omega$, using a $C^0$ tolerance smaller than both $\varepsilon$ and
half this positive lower bound.  Extend the perturbation by zero on the
sphere.  It yields the desired integral, preserves the envelope
inequalities, and preserves strictness and nonvanishing in $\Omega$.
Outside $\Omega$ the profile is unchanged, so it is strict throughout
$r\le3/10$ and has no critical points there except $r=0$.

At $r=0$, the perturbation vanishes on a neighborhood and
$\dd\phi_0=0$.  The connection term in the Hessian therefore vanishes,
and for every incoming metric
\[
 H_\phi=2(\dd x_1^2+\dd x_2^2)
 \qquad\text{at }r=0.
\]
This form is positive semidefinite of rank two.  It is positive on a
two-plane, and every two-plane contains a direction on which it is
strictly positive.  This last observation will allow transverse phase
directions to be chosen even at critical points.

Finally, shrink $\mathcal U$ if necessary so that the metric and
coordinate norms, and the coordinate volume densities on
$\overline\Omega$, are uniformly comparable.  Since $L_\phi\ge|a_\phi|$,
the uniform upper bound for $\dd V_g/\dd x$ proves
\eqref{eq:profile-coordinate-mass}.  Each coordinate function on
$\overline\Omega$ extends smoothly to the sphere using a cutoff in its
chart.  The extensions have uniformly bounded gradients for
$g\in\mathcal U$, so their two-coordinate projections have uniform
Lipschitz constants for the ambient Riemannian distance.
\end{proof}
